\section{Functional Requirements}
\begin{justify}
The functional requirements follow the IEEE SRS style by describing visible behavior expected from the software.
\end{justify}

\begin{table}[H]
\centering
\begin{tabular}{|M|L|}
\hline
\textbf{Requirement ID} & \textbf{Description} \\
\hline
FR-01 & The system shall allow the user to enter task name, expected time, and actual time. \\
\hline
FR-02 & The system shall validate that all required fields are present. \\
\hline
FR-03 & The system shall validate that expected time and actual time are positive numbers. \\
\hline
FR-04 & The system shall classify a task as Underestimation, Overestimation, or Accurate. \\
\hline
FR-05 & The system shall calculate the difference between expected and actual time. \\
\hline
FR-06 & The system shall calculate and return an accuracy percentage. \\
\hline
FR-07 & The system shall store task details and analysis results in MySQL. \\
\hline
FR-08 & The system shall provide APIs for registration, login, task list, dashboard, and suggestion. \\
\hline
\end{tabular}
\caption{Functional Requirements}
\end{table}

\section{Non Functional Requirements}
\begin{table}[H]
\centering
\begin{tabular}{|M|L|}
\hline
\textbf{Requirement ID} & \textbf{Description} \\
\hline
NFR-01 & The interface should be simple and understandable for first-time users. \\
\hline
NFR-02 & The backend should respond quickly for normal task submission. \\
\hline
NFR-03 & The database design should maintain task and analysis data consistently. \\
\hline
NFR-04 & The system should be maintainable by separating server logic, database connection, and SQL queries. \\
\hline
NFR-05 & The system should be extendable for future machine learning-based prediction. \\
\hline
\end{tabular}
\caption{Non Functional Requirements}
\end{table}

\section{Specific Requirements}
\begin{table}[H]
\centering
\begin{tabular}{|M|L|}
\hline
\textbf{Category} & \textbf{Requirement} \\
\hline
Hardware & Laptop or desktop computer with minimum 4 GB RAM and basic internet/browser support. \\
\hline
Software & Node.js, Express.js, MySQL Server, browser, code editor, and Overleaf/LaTeX for report preparation. \\
\hline
Frontend & HTML, CSS, and JavaScript. \\
\hline
Backend & Node.js, Express.js, CORS, mysql2 package. \\
\hline
Database & MySQL database named \texttt{task\_estimation\_system}. \\
\hline
\end{tabular}
\caption{Hardware and Software Requirements}
\end{table}

\section{UseCase Diagrams and Description}
\begin{justify}
The main actor is the user. The user can add a task, enter expected and actual time, view analysis, register, login, view tasks, view dashboard data, and get suggested time. The backend system validates input, stores data, and generates the estimation result.
\end{justify}

\begin{table}[H]
\centering
\begin{tabular}{|M|L|}
\hline
\textbf{Use Case} & \textbf{Description} \\
\hline
Add Task & User submits task name, expected time, and actual time. \\
\hline
Analyze Task & System compares expected and actual time and generates classification. \\
\hline
View Dashboard & User views previous tasks with analysis details. \\
\hline
Get Suggestion & System calculates average actual time for a matching task name. \\
\hline
\end{tabular}
\caption{Use Case Description}
\end{table}
